\documentclass[12pt]{article}
\usepackage{amsmath,amssymb}
\newcommand{\Dr}{\operatorname{Dr}}
\newcommand{\Trop}{\operatorname{Trop}}
\newcommand{\rk}{\operatorname{rk}}
\newcommand{\cl}{\operatorname{cl}}
\newcommand{\cL}{\mathcal L}
\begin{document}
\section*{Research notes}

Problem: Given a matroid $M$ and the relative Dressian $\Dr(M)$ of
tropical linear subspaces contained in $\Trop(M)$, suppose the lattice
of flats $\cL(M)$ has an order-reversing involution
$F\mapsto F^\perp$.  Need an order-reversing involution on $\Dr(M)$
extending
\[
F\longmapsto \Trop(M/F\oplus U_{0,F}).
\]

\subsection*{Conventions}

Use Brandt--Eur--Zhang conventions.  Let $\mu_M$ be the trivial rank
$r=\rk M$ valuated matroid supported on bases of $M$.  A rank $k$
element of $\Dr(M)$ is represented by a rank $k$ valuated matroid
$\nu$ with $\nu\preceq\mu_M$ (valuated quotient).  Inclusion is the
opposite order:
\[
\Trop(\theta)\subseteq\Trop(\nu)
\Longleftrightarrow
\theta\preceq\nu.
\]
Valuated matroids are projective: adding a common scalar to all
Plucker coordinates does not change the tropical linear space.

For valuated matroids $p,q$ of ranks $a<b$, the full incidence
relation is
\[
\mathcal I_{a,b}(p,q;S,T):\quad
\min_{i\in T\setminus S}\{p(S\cup i)+q(T\setminus i)\}
\text{ is attained at least twice},
\]
where $|S|=a-1$, $|T|=b+1$.  For consecutive ranks $d,d+1$ this is
equivalent, in the presence of the valuated-matroid axioms, to the
three-term incidence relations
\[
 \min\{p(Sx)+q(Syz),p(Sy)+q(Sxz),p(Sz)+q(Sxy)\}
\]
for $|S|=d-1$ and $x,y,z\notin S$.  This is the elementary quotient
criterion used in the current proof.

For complete sequences with trivial rank $0$ and rank $r$ endpoints,
the adjacent three-term relations alone characterize valuated complete
flags (Murota's $M^\natural$-concave/valuated generalized matroid
cryptomorphism).  Thus, after reversing a complete Higgs flag, it is
enough to verify adjacent three-term relations; individual Plucker
relations for the transformed slices follow from the complete-flag
criterion.


\subsection*{Structural consequence of the hypothesis}

A finite flat lattice is geometric and upper semimodular.  If it has
an anti-automorphism $x\mapsto x^\perp$, then ranks are reversed:
$\rk(x^\perp)=r-\rk(x)$.  Applying the anti-automorphism to
semimodularity gives the opposite inequality, hence
\[
\rk(x)+\rk(y)=\rk(x\wedge y)+\rk(x\vee y)
\]
for all $x,y$.  Therefore $\cL(M)$ is modular and
\[
(x\vee y)^\perp=x^\perp\wedge y^\perp,\qquad
(x\wedge y)^\perp=x^\perp\vee y^\perp.
\]

Loops/parallel elements: pass to the simplification.  The flat lattice
is canonically unchanged (up to identifying parallel classes), loops
are loops in every quotient, and flat constancy makes parallel
representatives have identical coordinates.  The formula for the
orthogonal transform is phrased only in terms of closures/flats, so it
pulls back from the simplification.

\subsection*{Flat constancy}

If $\nu\preceq\mu_M$ has rank $k$, then $\nu(B)$ depends only on
$\cl_M(B)$ for independent $k$-sets $B$.

Adjacent-basis proof for $0<k<r$: compare $A b$ and $A b'$ with
$\cl(A b)=\cl(A b')$.  Choose $C$ of size $r-k$ such that $A b C$ is
a basis of $M$; then $A b' C$ is also a basis.  In
$\mathcal I_{k,r}(\nu,\mu_M;A,A b b' C)$ only the $b,b'$ terms have
finite $\mu_M$ factor, since deleting an element of $C$ leaves the
dependent subset $A b b'$.  Thus $\nu(A b)=\nu(A b')$.  Endpoints:
$k=0$ trivial; $k=r$ gives the top element $\mu_M$ up to scalar.

Also use the quotient support property: if $\nu(A)<\infty$, then $A$
is independent in $M$.

Thus a rank $k$ quotient gives flat coordinates
\[
\bar\nu(X)=\nu(B)\quad\text{for }\rk X=k,\ \cl(B)=X,
\]
and $\nu(A)=\bar\nu(\cl A)$ for independent $A$, while dependent
$A$ has value $\infty$.

\subsection*{Candidate transform}

For a rank $k$ quotient $\nu$:
\[
\nu^{\perp_M}(A)=
\begin{cases}
\bar\nu((\cl A)^\perp),& A\text{ independent in }M,\ |A|=r-k,\\
\infty,&\text{otherwise}.
\end{cases}
\]
This is involutive in flat coordinates.  Adding a scalar to $\nu$
adds the same scalar to $\nu^{\perp_M}$.

\subsection*{Valuated Higgs factorization}

Key standard input: every valuated quotient $p\preceq q$ of rank
difference $b-a$ factors into elementary quotients
\[
p=h_a\preceq h_{a+1}\preceq\cdots\preceq h_b=q.
\]
For $q=\mu_M$, the upper Higgs lifts from a rank $a$ quotient $p$ can
be written on independent $j$-sets $B$ of $M$ as
\[
h_j(B)=\min\{p(A):A\subseteq B,\ |A|=a\},\qquad a\le j\le r,
\]
with value $\infty$ on dependent sets.  The rank $r$ lift is scalar
$\mu_M$; constancy over bases follows directly from the incidence
relations with $\mu_M$ by adjacent basis exchange.

Polyhedral intuition: a valuated quotient $p\preceq\mu_M$ induces a
regular subdivision of the quotient matroid into matroids that are
still quotients of $M$ (initial flags of a valuated flag matroid).
Taking ordinary Higgs lifts of each cell toward $M$ gives matroidal
subdivisions of the Higgs-lift polytopes.  This is the M-convex
independent-set viewpoint of Murota/Dress--Wenzel.  The current proof
uses this as a cited standard theorem; if challenged, provide a
separate appendix proving the needed factorization by M-convex
convolution or by regular subdivisions of Higgs lift polytopes.

\subsection*{Local elementary opposition}

Let $p\preceq q\preceq\mu_M$ with $\rk p=d$, $\rk q=d+1$.  The local task is to check the elementary three-term
relation for $q^\perp,p^\perp$; in a transformed complete flag the complete-flag criterion then promotes all adjacent checks to actual valuated quotients.  Fix $A$ of size $r-d-2$ and
$x,y,z\notin A$.  Let $U=\cl(A)$ and $W=\cl(Axyz)$, with
$\delta=\rk W-\rk U$.

If $A$ is dependent or $\delta\le1$, all relevant terms are $\infty$.

If $\delta=2$, set $G=W^\perp$ (rank $d$) and $H=U^\perp$ (rank
$d+2$).  All finite transformed terms share the same $p$-coordinate
$\bar p(G)$.  The $q$-coordinates are rank $d+1$ flats
$L_x=(\cl Ax)^\perp$, etc., in the rank-two interval $[G,H]$.  If
some $L$'s coincide, the corresponding finite terms are equal.  If
they are three distinct points, use the quotient relation
$q\preceq\mu_M$: choose a basis $S$ of $G$, representatives of
$L_x,L_y,L_z$ in $H$, and a complement $C$ so that deleting one of
the representatives from $S\cup\{\ell_x,\ell_y,\ell_z\}\cup C$ gives
a basis of $M$, while deleting from $C$ leaves a dependent line.  This
forces the minimum of $\bar q(L_x),\bar q(L_y),\bar q(L_z)$ to occur
twice.

If $\delta=3$, set $G=W^\perp$ (rank $d-1$) and $H=U^\perp$ (rank
$d+2$).  Define
\[
P_x=(\cl(Ayz))^\perp,\quad P_y=(\cl(Axz))^\perp,\quad
P_z=(\cl(Axy))^\perp
\]
and
\[
Q_x=(\cl(Ax))^\perp,\quad Q_y=(\cl(Ay))^\perp,\quad
Q_z=(\cl(Az))^\perp.
\]
In the rank-three modular interval $[G,H]$, the $P$'s are independent
atoms over $G$, and $Q_x=P_y\vee P_z$, etc.  Choose representatives
$a_x\in P_x\setminus G$, etc.  The original elementary incidence
relation for $p\preceq q$ on a basis $S$ of $G$ and
$a_x,a_y,a_z$ is exactly the transformed relation.

This avoids the previous need to dualize arbitrary subset-indexed
relations directly; arbitrary rank differences are handled by Higgs
factorization into elementary steps.

\subsection*{Order reversal and flat extension}

Given $\nu\preceq\mu_M$, put it into a complete factorization
$h_0\preceq\cdots\preceq h_r=\mu_M$ with $h_k=\nu$.  Elementary
opposition gives
\[
h_r^\perp\preceq h_{r-1}^\perp\preceq\cdots\preceq h_0^\perp,
\]
where $h_r^\perp$ is rank $0$ and $h_0^\perp=\mu_M$.  Hence
$\nu^\perp\preceq\mu_M$.  If $\theta\preceq\nu$, factor a complete
flag containing $\theta,\nu,\mu_M$; after opposition it contains
$\nu^\perp\preceq\theta^\perp$, which is order reversal in inclusion
of tropical linear spaces.

For an embedded flat $F$ of rank $f$, $q_F=M/F\oplus U_{0,F}$ has
finite rank $r-f$ flat coordinate at $Z$ iff $Z\vee F=\hat1$.  For an
independent $f$-set $A$ with $Y=\cl(A)$,
\[
\Phi(q_F)(A)<\infty
\Longleftrightarrow Y^\perp\vee F=\hat1
\Longleftrightarrow Y^\perp\wedge F=\hat0
\Longleftrightarrow Y\vee F^\perp=\hat1,
\]
using modularity and complementary ranks.  This is exactly the basis
condition for $q_{F^\perp}=M/F^\perp\oplus U_{0,F^\perp}$.

\subsection*{Computational sanity checks from round 1}

The attached compute-worker script checked the full subset-indexed
Plucker and incidence relations, not just interval summaries, for:
Boolean $U_{4,4}$, $U_{2,4}$, $U_{2,5}$, Fano $PG(2,2)$ with
dot-product polarity, $U_{2,3}\oplus U_{2,3}$ (including an exhaustive
finite rank-two search over $\{0,1,2\}$ up to normalization), and
random samples in $PG(2,3)$.  No counterexample to the flat-coordinate
opposition was found.  The tests also support factorization by the
explicit relative Higgs lifts
$h_j(B)=\min\{p(A):A\subseteq B\}$ in the small examples.

\end{document}
